% P2 poster: Design and Evaluation of a Drift-Aware Continuous Learning MLOps Framework for Financial Transaction Fraud Detection
% Generated by build_poster_latex.py. 48 x 36 in (121.92 x 91.44 cm), landscape. Compile with pdfLaTeX.
\documentclass[final]{beamer}
\usepackage[size=custom,width=121.92,height=91.44,scale=1.0]{beamerposter}
\usepackage[T1]{fontenc}
\usepackage[default]{lato}
\usepackage{graphicx,booktabs,array,colortbl}
\usepackage[most]{tcolorbox}

% ---------------------------------------------------------------- colours (from the BRAC template)
\definecolor{bracnavy}{HTML}{003C77}
\definecolor{pinkbox}{HTML}{FFF5F2}
\definecolor{greybox}{HTML}{F0F0F0}
\definecolor{ink}{HTML}{1A1A1A}
\definecolor{rule}{HTML}{555555}
\definecolor{rowalt}{HTML}{F3F6FA}
\setbeamercolor{normal text}{fg=ink,bg=white}
\setbeamertemplate{navigation symbols}{}
\setlength{\parskip}{0.45ex}
\setlength{\tabcolsep}{0.5em}

% ---------------------------------------------------------------- page geometry
\newlength{\sepwidth}\setlength{\sepwidth}{0.022\paperwidth}
\newlength{\colwidth}\setlength{\colwidth}{0.3\paperwidth}
\newcommand{\separatorcolumn}{\begin{column}{\sepwidth}\end{column}}

% ---------------------------------------------------------------- header and footer
\setbeamertemplate{headline}{%
  {\setlength{\fboxsep}{0pt}\colorbox{bracnavy}{\begin{minipage}[c][15.5cm][c]{\paperwidth}
    \centering
    \begin{minipage}[c]{0.78\paperwidth}\centering\color{white}
      {\fontsize{78}{94}\selectfont\bfseries Design and Evaluation of a Drift-Aware Continuous Learning MLOps Framework\\ for Financial Transaction Fraud Detection\par}
      \vspace{0.9cm}
      {\fontsize{40}{54}\selectfont Presented by Mashud Hasan (22299053), Nirnoy Biswas (24241185), Md. Sabbir Hossain (22299038),\\ Al Imran (231011035), Imtiaz Zaman Sami (23101551)\\
       Supervisor: Md. Sabbir Ahmed\quad|\quad Co-Supervisor: Naima Tahsin Nodi\par}
      \vspace{0.5cm}
      {\fontsize{30}{36}\selectfont BRAC University\par}
    \end{minipage}%
    \hspace{0.02\paperwidth}%
    \begin{minipage}[c]{0.1\paperwidth}\includegraphics[width=\linewidth]{figures/brac_logo.png}\end{minipage}
  \end{minipage}}}}
\setbeamertemplate{footline}{%
  {\setlength{\fboxsep}{0pt}\colorbox{bracnavy}{\begin{minipage}[c][2.5cm][c]{\paperwidth}\color{white}\fontsize{28}{34}\selectfont
    \hspace{2.5cm}https://www.bracu.ac.bd\hfill
    BRAC University | Kha 224 Pragati Sarani, Merul Badda, Dhaka 1212 | Tel: +88 09638464646\hfill
    \copyright\ 2026 Brac University. All rights reserved.\hspace{2.5cm}
  \end{minipage}}}}

% ---------------------------------------------------------------- building blocks
\newcommand{\heading}[1]{\vspace{1.0cm}{\centering\color{bracnavy}\fontsize{40}{48}\selectfont\bfseries #1\par}
  \vspace{0.25cm}{\color{rule}\rule{\linewidth}{2.5pt}}\par\vspace{0.4cm}}
\newtcolorbox{pinkbox}{colback=pinkbox,colframe=pinkbox,sharp corners,boxrule=0pt,
  left=0.8cm,right=0.8cm,top=0.6cm,bottom=0.6cm}
\newtcolorbox{greybox}{colback=greybox,colframe=greybox,sharp corners,boxrule=0pt,
  left=0.7cm,right=0.7cm,top=0.4cm,bottom=0.6cm}
\newcommand{\fig}[3]{\begin{center}\includegraphics[width=#2\linewidth]{figures/#1}\par
  \vspace{0.2cm}{\fontsize{21}{26}\selectfont\color{rule}#3\par}\end{center}\vspace{-0.4cm}}
\setbeamertemplate{itemize items}{\textbullet}
\setbeamerfont{itemize/enumerate body}{size=\normalsize}
\setbeamertemplate{enumerate items}[default]
\setbeamercolor{item}{fg=ink}
\setbeamercolor{enumerate item}{fg=ink}

\begin{document}
\begin{frame}[t]
\vspace{-0.5cm}
\begin{columns}[t]
\separatorcolumn

% ================================================================ column 1
\begin{column}{\colwidth}
\heading{Abstract}
Fraud models are usually trained once, but in production fraud patterns change
(\textbf{concept drift}) and true labels arrive weeks late (\textbf{label delay}). We build and
evaluate a \textbf{drift-aware continuous learning MLOps framework} that monitors a deployed
model on delayed labels, retrains it on a schedule with an automatically chosen training
window, and deploys a new model only if it passes a \textbf{promotion gate}. It is tested on
three public datasets with label delays of 0 to 60 days.

\heading{Problem Statement}
\begin{pinkbox}
Existing studies test fraud models on \textbf{random splits}, assume \textbf{immediate labels}
and trigger retraining with \textbf{drift detectors} on artificial drift. When should a real
fraud model be retrained, and how can a bad model be kept out of production?
\end{pinkbox}

\heading{Research Objectives}
\begin{enumerate}
  \item Measure how a deployed fraud model decays and whether drift measures detect it.
  \item Measure how label delay limits the benefit of retraining.
  \item Compare scheduled and drift-triggered retraining fairly.
  \item Build an MLOps framework with monitoring, retraining, window selection and a
        promotion gate, and test it with injected faults.
  \item Replicate on three public datasets.
\end{enumerate}

\heading{Literature Review}
\begin{itemize}
  \item \textbf{Dal Pozzolo et al.} [1]: fraud labels arrive late (verification latency),
        which limits how a model can be updated.
  \item \textbf{Menezes and Filho} [2]: fraud models on IEEE-CIS lose up to 40\% F1 under
        natural drift; no retraining tested.
  \item \textbf{Shakil et al.} [3]: drift-triggered retraining for MLOps, tested only on
        injected drift with immediate labels.
\end{itemize}
\textbf{Gap:} no fair test of retraining policies for fraud under real drift and label delay.

\heading{Data set}
{\fontsize{21}{27}\selectfont
\begin{tabular}{>{\raggedright\arraybackslash}p{0.2\linewidth}
                >{\raggedright\arraybackslash}p{0.24\linewidth}
                >{\raggedright\arraybackslash}p{0.24\linewidth}
                >{\raggedright\arraybackslash}p{0.22\linewidth}}
\rowcolor{bracnavy}\color{white} & \color{white}\textbf{IEEE-CIS} & \color{white}\textbf{Sparkov} & \color{white}\textbf{BAF}\\
Type & Real e-commerce & Simulated cards & Bank accounts\\
\rowcolor{rowalt} Records & 590,540 & 1,801,969 & 1,000,000\\
Time span & 182 days & 720 days & 8 months\\
\rowcolor{rowalt} Fraud rate & 3.50\% & 0.53\% & 1.10\%\\
Features & 431 & 16 & 29\\
\arrayrulecolor{bracnavy}\specialrule{2pt}{0pt}{0pt}
\end{tabular}\par}
\vspace{0.6cm}
\textbf{Preprocessing}
\begin{itemize}
  \item \textbf{IEEE-CIS:} transaction and identity tables joined on TransactionID; text
        columns as categories; missing values kept.
  \item \textbf{Sparkov:} personal fields removed; cut at 21 Dec 2020 (fraud stops); age,
        hour, distance and card-history features from \emph{earlier} transactions only.
  \item \textbf{BAF:} $-1$ codes set to missing; month not used as a feature.
  \item \textbf{All:} sorted by time, never shuffled; no oversampling (fraud class weighted).
\end{itemize}
\fig{preprocessing.png}{0.8}{Preprocessing pipeline from raw files to the time-ordered table
used by every experiment}
\end{column}

\separatorcolumn
% ================================================================ column 2
\begin{column}{\colwidth}
\heading{Methodology}
\fig{methodology_3ds.png}{0.82}{Workflow of the study}
\fig{model_architecture.png}{0.82}{Model architecture: LightGBM gradient-boosted decision trees}
\fig{architecture.png}{0.82}{Proposed framework: each week it monitors, decides, retrains,
gates and serves}
\vspace{0.6cm}
\begin{greybox}
\heading{Result Analysis}
\fig{results_decay.png}{0.92}{\textbf{Drift measures do not track performance:} IEEE-CIS
decays while drift stays low; Sparkov and BAF drift without decaying}
\end{greybox}
\end{column}

\separatorcolumn
% ================================================================ column 3
\begin{column}{\colwidth}
\vspace{1.0cm}
\begin{greybox}
\fig{results_delay.png}{1.0}{\textbf{Label delay erodes retraining:} on IEEE-CIS the gain
falls from +0.078 to +0.010 PR-AUC}
\fig{results_window.png}{1.0}{\textbf{Automatic window selection} matches the best fixed
window without hindsight}
\fig{results_gate.png}{0.9}{\textbf{The promotion gate} blocked every model trained on
corrupted data}
\end{greybox}

\heading{Conclusion \& Future Work}
\begin{itemize}
  \item \textbf{Drift indicators can miss real degradation under delayed labels}, while
        scheduled retraining was more reliable. Here, \emph{drift-aware} means tracking drift
        through delayed-label performance, not label-free alarms.
  \item With delayed labels, the schedule matched or beat the drift trigger even at an
        \textbf{equal retraining budget} (IEEE-CIS, BAF): not only from retraining more often.
  \item \textbf{Window selection} and a \textbf{promotion gate} make retraining adaptive and
        safe; results replicated on \textbf{two more timestamped datasets} (Sparkov, BAF).
  \item \textbf{API load test} (supervisor feedback): after a fix to the scoring path, about
        \textbf{1,300 requests/s}, p95 latency 82\,ms at 100 clients, identical scores.
        \textbf{Next:} equal-budget comparison on Sparkov.
  \item \textbf{Later:} faster investigator feedback, more model families, explanation
        monitoring and a live deployment.
\end{itemize}

\heading{References}
{\fontsize{18}{23}\selectfont
\begin{enumerate}
  \item[{[1]}] A. Dal Pozzolo et al., ``Credit card fraud detection: A realistic modeling and a
        novel learning strategy,'' \emph{IEEE Trans. Neural Netw. Learn. Syst.}, vol. 29,
        no. 8, 2018.
  \item[{[2]}] R. S. Menezes and R. H. Filho, ``Semantic and structural drift in financial
        knowledge graphs: A robustness analysis of GNN-based fraud detectors,'' in
        \emph{Proc. IEEE ICKG}, 2025.
  \item[{[3]}] M. N. P. Shakil et al., ``Feature drift-guided adaptive ML retraining: An MLOps
        approach for big data analytics,'' in \emph{Proc. IEEE BigData}, 2025.
\end{enumerate}\par}
\end{column}

\separatorcolumn
\end{columns}
\end{frame}
\end{document}
